\documentclass[11pt,letterpaper]{article}
\usepackage{graphicx} 
\usepackage{multirow}
\usepackage{enumitem}
\usepackage{amssymb}
\usepackage{amsmath}
\usepackage{amsthm}
\usepackage{xcolor}
\usepackage{array}
\usepackage{geometry}
  \geometry{
    left = 20mm,
    right = 20mm,
    top = 20mm,
    bottom = 30mm,
  }
\usepackage{fancyhdr}
\renewcommand{\headrulewidth}{0pt}
\pagestyle{fancy}

\begin{document}
\pagenumbering{gobble}

\begin{center}
  \textbf{\large Kuis 1 Matematika Statistika bab 7} \\
  \vspace{0.3cm}
  \textbf{\large Waktu: Pk. 11.00 -- 12.00 Wib}
\end{center}

\vspace{0.2cm}
\noindent Jawaban ditulis dalam kertas HVS dan dikumpulkan di my classroom ITS dalam 1 file bentuk pdf dan dikumpulkan paling lambat pk 12.05 \\
\rule{\textwidth}{0.5pt}
\vspace{0.2cm}

\begin{enumerate}
  \item Perhatikan suatu sampel random $X$ berukuran $n$ yang mempunyai CDF
  \[ F(x) = (1 + e^{-x})^{-1} ; x \in \mathbb{R} \]
  Dapatkan limiting distribution dari $\omega_n = X_{n:n} - \ln n$ \hfill [30]
  \vspace{0.3cm}

  \item Misalkan $Y_i$ adalah berat bagasi penumpang ke-$i$ dari suatu pesawat terbang (dalam kg). Diasumsikan beratnya saling independent dengan pdf tiap bagasi adalah
  $f(y_i) = \theta B^{-\theta} y_i^{\theta-1} ; 0 < y_i < B$. Jika diketahui $n = 100, \theta = 3, \text{dan } B = 80$,
  \begin{enumerate}[label=\alph*)]
    \item Dapatkan $E(y_i)$ dan $\text{Var}(y_i)$ \hfill [10]
    \item Dapatkan $E\left( \sum_{i=1}^{100} y_i \right)$ dan $\text{Var}\left( \sum_{i=1}^{100} y_i \right)$ \hfill [10]
    \item Dapatkan Peluang total berat 100 bagasi paling sedikit 6025 kg \hfill [15]
  \end{enumerate}
  \vspace{0.3cm}

  \item Misalkan $n$ sampel random dari distribusi Poisson, $X_i \sim Poi(\lambda)$.
  \begin{enumerate}[label=\alph*)]
    \item Tunjukkan bahwa $\bar{X}_n$ konvergen stokastik ke $\lambda$ \hfill [20]
    \item Tunjukan bahwa $\bar{X}_n e^{-\bar{X}_n}$ konvergen stokastik ke $P(X=1) = \lambda e^{-\lambda}$. \hfill [15]
  \end{enumerate}
\end{enumerate}

\vspace{0.5cm}
\begin{center}
  \textbf{Selamat bekerja, bekerjalah dengan jujur} \\
  \vspace{0.2cm}
  \textbf{Setialah dalam hal hal kecil}
\end{center}

\newpage
\begin{center}
  \textbf{\large SOLUSI}
\end{center}

\begin{enumerate}
  \item 
  Diketahui sampel acak berukuran $n$ dengan CDF:
  \[ F(x) = (1 + e^{-x})^{-1} = \frac{1}{1 + e^{-x}}, \quad x \in \mathbb{R} \]
  Misalkan $X_{n:n} = \max(X_1, X_2, \cdots, X_n)$ adalah statistik orde ke-$n$.
  CDF dari $X_{n:n}$ adalah:
  \[ F_{X_{n:n}}(x) = [F(x)]^n = \left( \frac{1}{1 + e^{-x}} \right)^n \]
  Kita ingin mencari limiting distribution dari $\omega_n = X_{n:n} - \ln n$.
  Mula-mula, kita cari CDF dari $\omega_n$, misalkan $F_{\omega_n}(w)$:
  \begin{align*}
    F_{\omega_n}(w) &= P(\omega_n \le w) \\
    &= P(X_{n:n} - \ln n \le w) \\
    &= P(X_{n:n} \le w + \ln n) \\
    &= F_{X_{n:n}}(w + \ln n) \\
    &= \left( \frac{1}{1 + e^{-(w + \ln n)}} \right)^n
  \end{align*}
  Perhatikan bahwa $e^{-(w + \ln n)} = e^{-w} e^{-\ln n} = e^{-w} \cdot \frac{1}{n} = \frac{e^{-w}}{n}$.
  Substitusikan bentuk ini ke dalam persamaan sebelumnya:
  \[ F_{\omega_n}(w) = \left( \frac{1}{1 + \frac{e^{-w}}{n}} \right)^n = \left( 1 + \frac{e^{-w}}{n} \right)^{-n} \]
  Untuk mencari limiting distribution, kita hitung limit dari $F_{\omega_n}(w)$ saat $n \to \infty$. 
  Kita gunakan sifat limit: $\lim_{n \to \infty} \left(1 + \frac{x}{n}\right)^n = e^x$.
  Maka,
  \begin{align*}
    \lim_{n \to \infty} F_{\omega_n}(w) &= \lim_{n \to \infty} \left[ \left( 1 + \frac{e^{-w}}{n} \right)^n \right]^{-1} \\
    &= \left( e^{e^{-w}} \right)^{-1} \\
    &= e^{-e^{-w}}
  \end{align*}
  Jadi, limiting distribution dari $\omega_n$ adalah $F(w) = e^{-e^{-w}}$ untuk $w \in \mathbb{R}$ (yang merupakan CDF dari distribusi Gumbel).

  \item
  Diketahui $Y_i$ dengan pdf:
  \[ f(y_i) = \theta B^{-\theta} y_i^{\theta-1}, \quad 0 < y_i < B \]
  Dengan parameter $n = 100, \theta = 3, B = 80$.
  \begin{enumerate}[label=\alph*)]
    \item Dapatkan $E(y_i)$ dan $\text{Var}(y_i)$.
    \begin{align*}
      E(y_i) &= \int_{-\infty}^{\infty} y f(y) \,dy = \int_0^B y \left( \theta B^{-\theta} y^{\theta-1} \right) dy \\
      &= \theta B^{-\theta} \int_0^B y^\theta \,dy = \theta B^{-\theta} \left[ \frac{y^{\theta+1}}{\theta+1} \right]_0^B \\
      &= \theta B^{-\theta} \left( \frac{B^{\theta+1}}{\theta+1} \right) = \frac{\theta}{\theta+1} B
    \end{align*}
    Substitusi nilai $\theta = 3$ dan $B = 80$:
    \[ E(y_i) = \frac{3}{3+1} (80) = \frac{3}{4} (80) = 60 \]
    Untuk mencari variansi, cari $E(y_i^2)$ terlebih dahulu:
    \begin{align*}
      E(y_i^2) &= \int_0^B y^2 \left( \theta B^{-\theta} y^{\theta-1} \right) dy = \theta B^{-\theta} \int_0^B y^{\theta+1} \,dy \\
      &= \theta B^{-\theta} \left[ \frac{y^{\theta+2}}{\theta+2} \right]_0^B = \theta B^{-\theta} \left( \frac{B^{\theta+2}}{\theta+2} \right) = \frac{\theta}{\theta+2} B^2
    \end{align*}
    Substitusi nilai $\theta = 3$ dan $B = 80$:
    \[ E(y_i^2) = \frac{3}{3+2} (80^2) = \frac{3}{5} (6400) = 3840 \]
    Maka variansinya adalah:
    \begin{align*}
      \text{Var}(y_i) &= E(y_i^2) - (E(y_i))^2 \\
      &= 3840 - (60)^2 = 3840 - 3600 = 240
    \end{align*}

    \item Dapatkan $E\left( \sum_{i=1}^{100} y_i \right)$ dan $\text{Var}\left( \sum_{i=1}^{100} y_i \right)$.
    Berdasarkan sifat ekspektasi dan variansi untuk variabel yang saling independen:
    \begin{align*}
      E\left( \sum_{i=1}^{100} y_i \right) &= \sum_{i=1}^{100} E(y_i) = 100 \times 60 = 6000 \\
      \text{Var}\left( \sum_{i=1}^{100} y_i \right) &= \sum_{i=1}^{100} \text{Var}(y_i) = 100 \times 240 = 24000
    \end{align*}

    \item Dapatkan peluang total berat 100 bagasi paling sedikit 6025 kg.
    Misalkan $S = \sum_{i=1}^{100} y_i$ adalah total berat bagasi. 
    Karena $n=100$ (cukup besar), berdasarkan Teorema Limit Pusat (CLT), distribusi total $S$ mendekati distribusi Normal dengan rata-rata $\mu_S = 6000$ dan variansi $\sigma_S^2 = 24000$.
    Sehingga $S \sim N(6000, 24000)$ secara asimtotik.
    Peluang $S \ge 6025$ adalah:
    \begin{align*}
      P(S \ge 6025) &= P\left( \frac{S - 6000}{\sqrt{24000}} \ge \frac{6025 - 6000}{\sqrt{24000}} \right) \\
      &\approx P\left( Z \ge \frac{25}{\sqrt{24000}} \right) \\
      &= P\left( Z \ge \frac{25}{154.919} \right) \\
      &\approx P(Z \ge 0.161)
    \end{align*}
    Dari tabel distribusi Normal Standar, $P(Z \le 0.16) \approx 0.5636$.
    Maka,
    \[ P(Z \ge 0.161) = 1 - P(Z \le 0.161) \approx 1 - 0.5636 = 0.4364 \]
    Jadi, peluang total berat 100 bagasi paling sedikit 6025 kg adalah sekitar 0.4364 atau 43.64\%.
  \end{enumerate}

  \item
  Misalkan $n$ sampel random dari distribusi Poisson, $X_i \sim Poi(\lambda)$.
  Untuk distribusi Poisson, $E(X_i) = \lambda$ dan $\text{Var}(X_i) = \lambda$.
  \begin{enumerate}[label=\alph*)]
    \item Tunjukkan bahwa $\bar{X}_n$ konvergen stokastik ke $\lambda$. \\
    Kita akan menunjukkan konvergensi dalam probabilitas (konvergen stokastik), yang dilambangkan dengan $\bar{X}_n \xrightarrow{P} \lambda$.
    Kita bisa menggunakan Ketaksamaan Chebyshev:
    Untuk setiap $\epsilon > 0$,
    \[ P(|\bar{X}_n - E(\bar{X}_n)| \ge \epsilon) \le \frac{\text{Var}(\bar{X}_n)}{\epsilon^2} \]
    Kita tahu $E(\bar{X}_n) = E(X_i) = \lambda$ dan $\text{Var}(\bar{X}_n) = \frac{\text{Var}(X_i)}{n} = \frac{\lambda}{n}$.
    Maka,
    \[ P(|\bar{X}_n - \lambda| \ge \epsilon) \le \frac{\lambda/n}{\epsilon^2} = \frac{\lambda}{n \epsilon^2} \]
    Jika $n \to \infty$, maka limit dari batas atas tersebut adalah 0:
    \[ \lim_{n \to \infty} P(|\bar{X}_n - \lambda| \ge \epsilon) \le \lim_{n \to \infty} \frac{\lambda}{n \epsilon^2} = 0 \]
    Karena peluang tidak mungkin negatif, maka $\lim_{n \to \infty} P(|\bar{X}_n - \lambda| \ge \epsilon) = 0$.
    Menurut definisi Hukum Kuat/Lemah Bilangan Besar (WLLN), hal ini menunjukkan bahwa rata-rata sampel $\bar{X}_n$ konvergen stokastik (konvergen dalam probabilitas) ke $\lambda$, yaitu $\bar{X}_n \xrightarrow{P} \lambda$.

    \item Tunjukkan bahwa $\bar{X}_n e^{-\bar{X}_n}$ konvergen stokastik ke $P(X=1) = \lambda e^{-\lambda}$. \\
    Berdasarkan Teorema Pemetaan Kontinu (Continuous Mapping Theorem / Teorema Slutsky), jika barisan variabel acak $Z_n \xrightarrow{P} c$ (konvergen stokastik ke konstanta $c$) dan $g$ adalah fungsi yang kontinu di $c$, maka $g(Z_n) \xrightarrow{P} g(c)$. \\
    Misalkan fungsi $g(x) = x e^{-x}$. Fungsi $g(x)$ merupakan fungsi yang kontinu untuk seluruh nilai $x$ riil.
    Dari pembuktian bagian (a), kita sudah mendapatkan bahwa $Z_n = \bar{X}_n \xrightarrow{P} \lambda$.
    Sehingga, berdasarkan Teorema Pemetaan Kontinu:
    \[ g(\bar{X}_n) \xrightarrow{P} g(\lambda) \]
    \[ \bar{X}_n e^{-\bar{X}_n} \xrightarrow{P} \lambda e^{-\lambda} \]
    Di sisi lain, untuk distribusi Poisson $X \sim Poi(\lambda)$, fungsi probabilitas massanya adalah $P(X=x) = \frac{\lambda^x e^{-\lambda}}{x!}$.
    Untuk $X = 1$, kita dapatkan:
    \[ P(X=1) = \frac{\lambda^1 e^{-\lambda}}{1!} = \lambda e^{-\lambda} \]
    Dengan demikian, terbukti bahwa $\bar{X}_n e^{-\bar{X}_n}$ konvergen stokastik ke $P(X=1) = \lambda e^{-\lambda}$.
  \end{enumerate}
\end{enumerate}

\end{document}
